\section*{Bab \ref{ch:b1:derivative} --- Pendiferensialan}

\begin{solution}{exo:b1:derivative:1}
$x^x = \eu^{x\ln x}$ pada $\intoo{0}{+\infty}$: dengan \omterm{def:b1:derivative:def}{turunan} $(\ln x +
1)\,x^x$.

$\ln(x + \sqrt{x^2+1})$ pada $\R$ (karena argumennya selalu $> 0$):
dengan \omterm{def:b1:derivative:def}{turunan} $\frac{1}{\sqrt{x^2+1}}$ (yang dihitung pada
\cref{prop:b1:functions:invhyp} --- karena ia
$\operatorname{arsinh}$).

$\arctan\frac1x$ pada $\R^*$: dengan \omterm{def:b1:derivative:def}{turunan} $\frac{-1/x^2}{1 + 1/x^2} =
\frac{-1}{1 + x^2}$ (yang selaras dengan
\cref{prop:b1:functions:arcidentities}~(2): karena fungsinya adalah
$\pm\frac\pi2 - \arctan x$ pada masing-masing setengah garisnya).

$\sqrt{1 + \eu^{2x}}$ pada $\R$: dengan \omterm{def:b1:derivative:def}{turunan}
$\frac{\eu^{2x}}{\sqrt{1 + \eu^{2x}}}$.
\end{solution}

\begin{solution}{exo:b1:derivative:2}
Di $0$: $\bigl|\frac{f(h) - 0}{h}\bigr| = \abs{h \sin\frac1h} \leq
\abs h \to 0$, sehingga $f'(0) = 0$. Untuk $x \neq 0$, aturan yang lazim memberikan
$f'(x) = 2x\sin\frac1x - \cos\frac1x$. Sepanjang $x_n = \frac{1}{2\pi
n}$: $f'(x_n) = 0 - 1 \to -1$; sedangkan sepanjang $y_n = \frac{1}{(2n+1)\pi}$:
$f'(y_n) = 0 + 1 \to 1$. Jadi dua barisan yang menuju $0$ dengan limit
$f'$ yang berbeda: sehingga tak ada limitnya (\cref{thm:b1:continuity:seqchar}), jadi $f'$
tak \omterm{def:b1:continuity:continuous}{kontinu} di $0$ dan $f$ \omterm{def:b1:derivative:def}{dapat diturunkan} tanpa menjadi
$C^1$.
\end{solution}

\begin{solution}{exo:b1:derivative:3}
$\ln(1+x) < x$ untuk $x > 0$: lewat ketaksamaan garis singgung kecekungannya di $0$
(yang tegas jauh dari titik sentuhnya karena $\ln$ cekung tegas;
atau terapkanlah teorema nilai rata-rata: $\ln(1+x) = \frac{x}{1+c}$ untuk suatu
$c \in \intoo{0}{x}$, dan $\frac{x}{1+c} < x$). Kesamaan nilai rata-rata
yang sama memberikan batas bawahnya: $\frac{x}{1+c} > \frac{x}{1+x}$.

Akibatnya: dengan $x/n$ menggantikan $x$,
\[
\frac{x/n}{1 + x/n} < \ln\Bigl(1 + \frac xn\Bigr) < \frac xn
\quad\implies\quad
\frac{x}{1 + x/n} < n \ln\Bigl(1 + \frac xn\Bigr) < x .
\]
Ruas kirinya menuju $x$: sehingga menurut apitannya, $n\ln(1 + \frac xn)
\to x$, lalu menurut \omterm{def:b1:continuity:continuous}{kekontinuan} $\exp$, $\bigl(1 + \frac xn\bigr)^n =
\eu^{n\ln(1 + x/n)} \to \eu^x$.
\end{solution}

\begin{solution}{exo:b1:derivative:4}
Misalkan $x_1 < x_2 < \dots < x_k$ akar $P$ yang berbeda. Pada setiap
$\intcc{x_i}{x_{i+1}}$, Rolle (\cref{thm:b1:derivative:rolle})
menghasilkan $c_i \in \intoo{x_i}{x_{i+1}}$ dengan $P'(c_i) = 0$: jadi itu
$k - 1$ akar $P'$, yang berbeda karena \omterm{prop:b1:reals:intervals}{selang} \omterm{def:b1:topology:open}{terbukanya}
lepas --- dan terselang-seling menurut konstruksinya.

Jika $P$ (yang berderajat $n$) berakar real semuanya, tulislah akarnya dengan
multiplisitas $m_1 + \dots + m_k = n$. Setiap akar bermultiplisitas
$m_i \geq 2$ merupakan akar $P'$ bermultiplisitas $m_i - 1$
(\cref{prop:b1:poly:multiplicity}), yang menyumbang $\sum (m_i - 1) = n
- k$; lalu Rolle menyumbang $k - 1$ lagi, yang semuanya berbeda dari itu. Jadi totalnya
$\geq n - 1 = \deg P'$: sehingga semua akar $P'$ bersifat real.
\end{solution}

\begin{solution}{exo:b1:derivative:5}
Tetapkanlah $\varepsilon > 0$ dan $A$ dengan $\abs{f'(t) - \ell} \leq
\varepsilon$ untuk $t \geq A$. Untuk $x > A$, teorema nilai rata-rata pada
$\intcc{A}{x}$ memberikan $c \in \intoo{A}{x}$ dengan
\[
f(x) = f(A) + f'(c)(x - A),
\qquad\text{sehingga}\qquad
\Bigl|\frac{f(x)}{x} - \ell\Bigr|
\leq \frac{\abs{f(A)} + \abs\ell A}{x} + \abs{f'(c) - \ell}
\cdot\frac{x - A}{x} \leq \frac{C_A}{x} + \varepsilon .
\]
Untuk $x$ yang besar, $\frac{C_A}{x} \leq \varepsilon$: sehingga
$\frac{f(x)}{x} \to \ell$.

Ya: karena $f(x+1) - f(x) = f'(c_x)$ dengan $c_x \in \intoo{x}{x+1}$ (menurut teorema
nilai rata-rata pada $\intcc{x}{x+1}$), dan $c_x \to +\infty$, sehingga
$f(x+1) - f(x) \to \ell$.
\end{solution}

\begin{solution}{exo:b1:derivative:6}
Lewat induksi pada $n$. Untuk $n = 1$: Rolle. Jika klaimnya berlaku untuk $n - 1$:
karena $f$ lenyap di $n+1$ titik, maka menurut Rolle yang diterapkan pada $n$ celahnya,
$f'$ lenyap di $n$ titik yang berbeda; lalu hipotesis induksinya
yang diterapkan pada $f'$ (yang \omterm{def:b1:derivative:def}{dapat diturunkan} $n-1$ kali, dengan $n$ nol) membuat
$(f')^{(n-1)} = f^{(n)}$ lenyap di suatu tempat.

Penerapannya: jika $P$ berderajat $\leq n$ lenyap di $n+1$ titik, maka
$P^{(n)}$, yaitu konstanta yang sama dengan $n!$ dikali koefisien utamanya,
lenyap: sehingga koefisien utamanya $0$, lalu kita menyimpulkannya lewat
induksi menurun (atau secara langsung: bahwa semua koefisiennya lenyap).
\end{solution}

\begin{solution}{exo:b1:derivative:7}
Menurut teorema nilai rata-rata pada $\intcc{a}{x_0}$ dan pada
$\intcc{x_0}{b}$:
\[
f'(c_1) = \frac{f(x_0) - f(a)}{x_0 - a} = \frac{f(x_0)}{x_0 - a} > 0,
\qquad
f'(c_2) = \frac{f(b) - f(x_0)}{b - x_0} = \frac{-f(x_0)}{b - x_0} < 0,
\]
dengan $c_1 < x_0 < c_2$. Lalu teorema nilai rata-rata yang diterapkan pada $f'$
pada $\intcc{c_1}{c_2}$ memberikan $c$ dengan
\[
f''(c) = \frac{f'(c_2) - f'(c_1)}{c_2 - c_1} < 0 . \qedhere
\]
\end{solution}

\begin{solution}{exo:b1:derivative:8}
$f'(x) = \frac{1 - \ln x}{x^2}$: jadi $f$ naik pada $\intoc{0}{\eu}$,
lalu turun pada $\intco{\eu}{+\infty}$, dengan maksimum $f(\eu) =
\frac1\eu$; dan limitnya $-\infty$ di $0^+$ serta $0$ di $+\infty$ (menurut
\omterm{prop:b1:functions:powerrules}{perbandingan pertumbuhannya}).

Untuk $\eu \leq a < b$: karena $f$ turun tegas di sana, maka
$\frac{\ln a}{a} > \frac{\ln b}{b}$, yakni $b \ln a > a \ln b$,
yakni $a^b > b^a$.

Dengan $a = \eu < b = \pi$: maka $\eu^\pi > \pi^\eu$.

Adapun pasangan yang kecil: $(2, 3)$: $2^3 = 8 < 9 = 3^2$ --- yang terbalik! Alasannya:
karena $2 < \eu$, dan pada $\intoo{0}{\eu}$ fungsi $f$ bersifat
\emph{naik}, sehingga perbandingannya membalik ketika kedua bilangannya duduk
di bawah $\eu$, dan tak teramalkan ketika melintasi $\eu$ (adapun $f(2) = f(4)$
menjelaskan seri $2^4 = 4^2 = 16$).
\end{solution}

\begin{solution}{exo:b1:derivative:9}
\emph{Jensen bagi $\ln$, lewat induksi pada $n$.} Klaimnya: bahwa untuk
$x_i$ yang positif dan bobot $\lambda_i > 0$ dengan $\sum \lambda_i = 1$:
$\ln\bigl(\sum \lambda_i x_i\bigr) \geq \sum \lambda_i \ln x_i$. Untuk
$n = 2$ inilah kecekungannya. Langkahnya: dengan $\Lambda = \lambda_1 + \dots +
\lambda_{n-1} = 1 - \lambda_n$ dan $y = \sum_{i<n}
\frac{\lambda_i}{\Lambda} x_i$,
\[
\ln\Bigl(\sum_{i \leq n} \lambda_i x_i\Bigr)
= \ln\bigl(\Lambda y + \lambda_n x_n\bigr)
\geq \Lambda \ln y + \lambda_n \ln x_n
\geq \Lambda \sum_{i<n} \frac{\lambda_i}{\Lambda}\ln x_i
+ \lambda_n \ln x_n,
\]
dengan memakai kecekungannya ($n = 2$) lalu hipotesis induksinya.

Dengan $\lambda_i = \frac 1n$ dan $x_i = a_i$: $\ln\frac{\sum a_i}{n}
\geq \frac 1n \sum \ln a_i = \ln\sqrt[n]{a_1\cdots a_n}$;
lalu eksponensialkanlah. Adapun kesamaannya: karena $\ln$ cekung \emph{tegas} ($\ln'' < 0$),
sehingga kesamaan pada setiap langkahnya memaksa titik yang dirata-ratakan berimpit ---
yakni semua $a_i$ sama; dan jika semuanya sama, kesamaannya jelas.
\end{solution}

\begin{solution}{exo:b1:derivative:10}
Misalkan $g(x) = f(x) - vx$: yang \omterm{def:b1:derivative:def}{dapat diturunkan}, dengan $g'(a) = f'(a) - v < 0$
dan $g'(b) = f'(b) - v > 0$. Menurut teorema nilai ekstrem, $g$
mencapai minimumnya pada $\intcc{a}{b}$ di suatu $c$. Dan ia tak di $a$:
karena $g'(a) < 0$, sehingga titik tepat di kanan $a$ mempunyai $g < g(a)$. Ia pun
tak di $b$: karena $g'(b) > 0$, sehingga titik tepat di kiri $b$ mempunyai $g <
g(b)$. Jadi $c$ bersifat \omterm{def:b1:topology:closure}{interior}, lalu \cref{prop:b1:derivative:fermat}
memberikan $g'(c) = 0$, yakni $f'(c) = v$.
\end{solution}

\begin{solution}{exo:b1:derivative:11}
\emph{Ketunggalannya:} dua titik tetap $\ell \neq \ell'$ akan memberikan
$\abs{\ell - \ell'} = \abs{f(\ell) - f(\ell')} \leq k\abs{\ell -
\ell'} < \abs{\ell - \ell'}$, yang mustahil.

\emph{Keberadaannya:} fungsi $g(x) = f(x) - x$ memenuhi, menurut ketaksamaan
nilai rata-rata, $f(x) \leq f(0) + k\abs x$; sehingga untuk $x \geq
\frac{\abs{f(0)}}{1 - k}$, $g(x) \leq f(0) + kx - x \leq 0$, dan
secara simetris $g(-x) \geq 0$ untuk $x$ yang besar. Lalu teorema nilai
antara memberikan sebuah nol $\ell$ bagi $g$: yaitu titik tetapnya.

\emph{Kekonvergenannya:} lewat ketaksamaan nilai rata-rata lagi:
\[
\abs{u_{n+1} - \ell} = \abs{f(u_n) - f(\ell)} \leq k\abs{u_n - \ell},
\]
sehingga lewat induksi $\abs{u_n - \ell} \leq k^n \abs{u_0 - \ell} \to 0$.
\end{solution}

\begin{solution}{exo:b1:derivative:12}
\begin{enumerate}
  \item Jika $g(b) = g(a)$, maka Rolle akan memberikan sebuah nol \omterm{def:b1:topology:closure}{interior} bagi
        $g'$: yang tersingkirkan. Tetapkanlah $\lambda = \frac{f(b) - f(a)}{g(b) -
        g(a)}$ dan $h = f - \lambda g$: maka $h$ \omterm{def:b1:continuity:continuous}{kontinu} pada
        $\intcc{a}{b}$, \omterm{def:b1:derivative:def}{dapat diturunkan} di dalamnya, dan $h(b) - h(a)
        = f(b) - f(a) - \lambda(g(b) - g(a)) = 0$. Lalu Rolle
        menyediakan $c$ dengan $h'(c) = 0$, yakni $f'(c) =
        \lambda\,g'(c)$; lalu bagilah dengan $g'(c) \neq 0$.
  \item Untuk $x > a$ yang dekat dengan $a$, bagian (1) pada $\intcc{a}{x}$
        (yang di situ $g' \neq 0$) memberikan $g(x) \neq 0$ dan $c_x \in
        \intoo{a}{x}$ dengan
        \[
        \frac{f(x)}{g(x)} = \frac{f(x) - f(a)}{g(x) - g(a)}
        = \frac{f'(c_x)}{g'(c_x)} .
        \]
        Ketika $x \to a^+$, $c_x \to a^+$ (lewat apitan), sehingga ruas
        kanannya menuju $\ell$: jadi $\frac{f}{g} \to \ell$.
  \item Di sini $\frac{f(x)}{g(x)} = x\sin\frac1x \to 0$, sedangkan
        $\frac{f'(x)}{g'(x)} = 2x\sin\frac1x - \cos\frac1x$ tak
        berlimit di $0$ (\cref{exo:b1:derivative:2}):
        jadi aturan l'Hospital memindahkan informasi hanya dari
        $\frac{f'}{g'}$ ke $\frac fg$, dan tak pernah sebaliknya.
\end{enumerate}
\end{solution}

\begin{solution}{pb:b1:derivative:1}
\textbf{1.} Di sini $\bigl|\frac ab - \frac pq\bigr| = \frac{\abs{aq -
bp}}{bq}$, dan $aq - bp$ merupakan bilangan bulat taknol ketika pecahannya
berbeda: sehingga jaraknya $\geq \frac{1}{bq}$. Jadi tak ada bilangan rasional
selain $x$ sendiri yang masuk ke \omterm{prop:b1:reals:intervals}{selang} berlubang berjari-jari
$\frac{1}{bq}$ di sekitar $x = \frac ab$.

\textbf{2.} Jika $\bigl|\sqrt2 - \frac pq\bigr| \geq 1 \geq
\frac{1}{4q^2}$, selesai. Kalau tidak, $\frac pq \in \intoo{\sqrt2 -
1}{\sqrt2 + 1}$, sehingga $0 < \sqrt2 + \frac pq < 2\sqrt2 + 1 < 4$.
Dan karena $\sqrt 2 \notin \Q$, maka $p^2 - 2q^2$ bilangan bulat taknol, sehingga
\[
\Bigl|\sqrt2 - \frac pq\Bigr|
= \frac{\abs{2q^2 - p^2}}{q^2\,\bigl(\sqrt2 + \frac pq\bigr)}
\geq \frac{1}{4q^2} .
\]

\textbf{3.} Di sini $(p + 2q)^2 - 2(p + q)^2 = -(p^2 - 2q^2)$: sehingga nilai
$\pm1$ merambat. Dari $(1,1)$:
\[
(3,2),\ (7,5),\ (17,12),\ (41,29),\ (99,70),
\]
dengan $p^2 - 2q^2$ berselang-seling $-1, +1, \dots$
Untuk pasangan ini, $\frac pq \geq 1$, sehingga $\sqrt2 + \frac pq > 2$ dan
\[
\Bigl|\sqrt2 - \frac pq\Bigr| =
\frac{1}{q^2(\sqrt2 + p/q)} < \frac{1}{2q^2} ,
\]
dengan $q \to \infty$: jadi hampiran berorde $2$ yang tak hingga banyaknya.
Digabung dengan pertanyaan 2, eksponen $2$ itu persis bagi $\sqrt 2$.

\textbf{4.} Sebanyak $N + 1$ bilangan $kx - \lfloor kx\rfloor$ (dengan $0
\leq k \leq N$) terletak pada $N$ kotak $\intco{\frac
jN}{\frac{j+1}{N}}$: sehingga dua di antaranya berbagi satu kotak
(\cref{cor:b1:counting:pigeonhole}), katakanlah untuk $i < j$. Dengan $q =
j - i \leq N$ dan $p = \lfloor jx\rfloor - \lfloor ix\rfloor$:
$\abs{qx - p} < \frac1N$, sehingga $\bigl|x - \frac pq\bigr| <
\frac{1}{Nq} \leq \frac{1}{q^2}$. Lalu membiarkan $N \to \infty$: karena
$x$ irasional, setiap pecahan yang tetap berjarak positif dari
$x$, sedangkan $\frac{1}{Nq} \leq \frac 1N \to 0$ memaksa pecahan
yang baru bermunculan: jadi ada $\frac pq$ berbeda yang tak hingga banyaknya dengan
$\bigl|x - \frac pq\bigr| < \frac{1}{q^2}$.

\textbf{5.} Mulailah dari sebarang $P_0$ bulat yang taknol dengan $P_0(x) =
0$. Jika $P_0$ berakar rasional $\frac ab$, maka teorema faktor
(\cref{thm:b1:poly:factor}) menulis $P_0 = \bigl(X - \frac
ab\bigr)Q$ dengan $Q \in \Q[X]$;
lalu karena $x \neq \frac ab$ ($x$ irasional), $Q(x) = 0$, dan
membersihkan penyebutnya memberikan \omterm{def:b1:poly:def}{polinomial} \emph{bulat} taknol
yang berderajat lebih kecil yang lenyap di $x$. Derajatnya turun pada setiap
langkahnya, sehingga prosesnya berhenti: jadi kita mencapai $P \in \Z[X]$, $P(x) = 0$,
tanpa akar rasional, berderajat $d$. Jika $d \leq 1$, maka $P =
uX + v$ akan membuat $x = -\frac vu$ rasional: sehingga $d \geq 2$.

\textbf{6.} Di sini $q^d\,P\bigl(\frac pq\bigr) = a_d p^d + a_{d-1}
p^{d-1} q + \dots + a_0 q^d$ merupakan bilangan bulat, dan ia taknol
karena $P$ tak berakar rasional: sehingga $\bigl|P\bigl(\frac
pq\bigr)\bigr| \geq q^{-d}$.

\textbf{7.} Perhatikan $M > 0$: karena $P'$ \omterm{def:b1:poly:def}{polinomial} taknol ($d \geq
2$), sehingga ia tak dapat lenyap secara identik pada
$\intcc{x-1}{x+1}$. Jika $\bigl|x - \frac pq\bigr| > 1$, maka ia
melampaui $\frac{C}{q^d}$ secara sepele. Kalau tidak, $\frac pq \in
\intcc{x-1}{x+1}$ dan teorema nilai rata-rata
(\cref{thm:b1:derivative:mvt}) memberikan $c$ di antara $x$ dan
$\frac pq$ dengan
\[
\Bigl|P\Bigl(\frac pq\Bigr)\Bigr|
= \Bigl|P\Bigl(\frac pq\Bigr) - P(x)\Bigr|
= \abs{P'(c)}\,\Bigl|x - \frac pq\Bigr|
\leq M\,\Bigl|x - \frac pq\Bigr| ,
\]
sehingga digabung dengan pertanyaan 6: $\bigl|x - \frac pq\bigr| \geq
\frac{1}{Mq^d} \geq \frac{C}{q^d}$.

\textbf{8.} Andaikan $x = \frac ab$ bersifat Liouville. Pilihlah $n$ dengan
$2^{n-1} > b$ beserta $\frac pq$ yang bersesuaian, dengan $q \geq 2$:
\[
0 < \Bigl|x - \frac pq\Bigr| < \frac{1}{q^n}
= \frac{1}{q^{n-1}\,q} \leq \frac{1}{2^{n-1} q} <
\frac{1}{bq} ,
\]
yang bertentangan dengan pertanyaan 1. Jadi bilangan Liouville bersifat irasional.

\textbf{9.} Seandainya $x$ yang Liouville itu aljabar: maka ia irasional
(pertanyaan 8), sehingga pertanyaan 5--7 menyediakan $d \geq 2$ dan $C > 0$
dengan $\bigl|x - \frac pq\bigr| \geq \frac{C}{q^d}$ selalu. Untuk
setiap $n$, penghampir Liouvillenya memberikan $\frac{C}{q^d} <
q^{-n}$, yakni $C < q^{d-n} \leq 2^{d-n}$ (karena $q \geq 2$). Lalu untuk
$n$ yang besar, $2^{d-n} < C$: yang bertentangan. Jadi bilangan Liouville bersifat
transenden.

\textbf{10.} Angka satunya pada posisi $1, 2, 6, 24$; sedangkan semua angka lain
di antara $25$ yang pertama lenyap:
\[
L = 0.1100010000\,0000000000\,00010\dots
\]

\textbf{11.} Untuk $m > k$, posisi $n!$ dengan $n > k$ merupakan
bilangan bulat berbeda yang $\geq (k+1)!$, sehingga jumlah geometri yang hingga
memberikan
\[
t_m - t_k = \sum_{n=k+1}^{m} 10^{-n!}
\leq \sum_{j = (k+1)!}^{m!} 10^{-j}
< 10^{-(k+1)!}\,\frac{1}{1 - \frac1{10}}
= \frac{10}{9}\,10^{-(k+1)!} ;
\]
lalu mengambil \omterm{def:b1:reals:bounds}{supremum} atas $m$: $L - t_k \leq
\frac{10}{9}10^{-(k+1)!} < 2\cdot10^{-(k+1)!}$. Adapun batas bawahnya: $L
\geq t_{k+1} = t_k + 10^{-(k+1)!}$.

\textbf{12.} Di sini $p_k = 10^{k!}\,t_k \in \N$, $q_k = 10^{k!}$, dan
$(k+1)! = (k+1)\,k!$ memberikan $10^{-(k+1)!} = q_k^{-(k+1)}$: sehingga
pertanyaan 11 berbunyi
\[
0 < L - \frac{p_k}{q_k} < \frac{2}{q_k^{\,k+1}} .
\]
Diberikan $n$: untuk $k \geq n$, $2\,q_k^{-(k+1)} \leq q_k^{-n}$
(memang $q_k^{\,k+1-n} \geq q_k \geq 10 > 2$), dan $q_k \geq 2$:
sehingga definisi pertanyaan 8 terpenuhi. Jadi $L$ merupakan bilangan Liouville.

\textbf{13.} Menurut pertanyaan 9, $L$ bersifat transenden --- yaitu bilangan
pertama dalam sejarah yang terbukti transenden (Liouville, 1844).
Adapun periksa silang angkanya: untaiannya memuat angka satu yang tak hingga banyaknya dengan
jarak berurutan $(k+1)! - k! = k\cdot k! \to \infty$, sehingga ia
tak periodik pada akhirnya, jadi $L \notin \Q$ menurut kriteria
keperiodikan pada \cref{pb:b1:reals:1} (pertanyaan 18) --- yang selaras.

\textbf{14.} Dengan angka $d_n \in \intint{1}{9}$ pada posisi
faktorialnya: batas ekor pertanyaan 11 terskala paling banyak
$9$: $0 < L' - t'_k \leq 9\cdot\frac{10}{9}\,10^{-(k+1)!} =
10\,q_k^{-(k+1)}$ (dengan kepositifannya karena angka pada posisi
$(k+1)!$ bersifat taknol). Untuk $k \geq n$: $10\,q_k^{-(k+1)} \leq
q_k^{-n}$ karena $q_k^{\,k+1-n} \geq 10$: jadi bilangan Liouville
lagi, sehingga transenden. Adapun nilainya berbeda berpasangan
untuk pemilihan angka yang berbeda (karena untaiannya sejati ---
sebab nolnya berlimpah --- dan untaian yang sejati menentukan nilainya,
\cref{pb:b1:reals:1}, pertanyaan 10). Diberikan sebarang daftar $k \mapsto
x_k$ atasnya, pilihlah angka faktorial ke-$k$ di
$\intint{1}{9}$ yang berbeda dari angka $x_k$: maka itu bilangan berbentuk
sama yang hilang dari daftarnya. Jadi \omterm{pb:b1:logic:1}{bilangan transenden} eksplisit yang
takterbilang banyaknya.

\textbf{15.} Mula-mula sebuah lema: \emph{bahwa jika $\bigl|x - \frac pq\bigr|
\geq \frac{C}{q^s}$ untuk setiap $\frac pq \neq x$, maka $x$ tak
terhampiri sampai orde $\mu > s$ mana pun.} Memang, adanya $\frac pq \neq x$
yang tak hingga banyaknya dengan $\bigl|x - \frac pq\bigr| <
\frac{c}{q^\mu}$ akan memaksa $\frac{C}{q^s} < \frac{c}{q^\mu}$,
yakni $q^{\mu - s} < \frac cC$: sehingga $q$-nya terbatas, dan
pecahan yang terbatas banyaknya terletak dalam jarak $1$ dari $x$ ---
jadi berhingga banyak calonnya, bukan tak hingga banyak. Sekarang
hierarkinya: bilangan rasional terhampiri sampai orde $1$ (karena $\frac
pq$ dengan $p = \lfloor qx\rfloor + 1$ memberikan galat $\leq \frac1q <
\frac2q$) dan tak sampai orde $\mu > 1$ mana pun (karena pertanyaan 1 memberikan
hipotesis lemanya dengan $s = 1$, $C = \frac1b$); adapun $\sqrt2$:
sampai orde $2$ (pertanyaan 3) dan tak lebih (pertanyaan 2 dan lemanya);
setiap bilangan irasional: sekurang-kurangnya $2$ (pertanyaan 4); yang aljabar berderajat
$d$: paling banyak $d$ (pertanyaan 7 dan lemanya); sedangkan bilangan Liouville:
setiap orde (menurut tampilan pertanyaan 12, dengan $c = 2$).

\textbf{16.} Dengan $r = \frac ab$: $\frac{p_k}{q_k} + \frac ab =
\frac{b p_k + a q_k}{b q_k} =: \frac{P_k}{Q_k}$, $Q_k = b q_k
\geq 2$, dan
\[
\Bigl|(L + r) - \frac{P_k}{Q_k}\Bigr| = L - \frac{p_k}{q_k}
< 2\,q_k^{-(k+1)} = 2\,b^{\,k+1} Q_k^{-(k+1)} .
\]
Diberikan $n$: untuk $k$ yang besar, $Q_k^{\,k+1-n} \geq Q_k = b\,10^{k!}
\geq 2\,b^{\,k+1}$ (karena faktorialnya menghancurkan pangkatnya), sehingga
galatnya $< Q_k^{-n}$: jadi $L + r$ bersifat Liouville. Dan karena $\Q$ \omterm{def:b1:topology:dense}{padat}
serta setiap $L + r$ bersifat transenden, maka \omterm{pb:b1:logic:1}{bilangan transenden}
bersifat \omterm{def:b1:topology:dense}{padat} di $\R$.

\textbf{17.} Untuk $h \geq 1$ ada berhingga banyak $P \in
\Z[X]$ dengan $\deg P + \sum_i \abs{a_i} \leq h$ (yaitu berderajat $\leq h$
dan setiap koefisiennya di $\intint{-h}{h}$: jadi paling banyak
$(2h+1)^{h+1}$). Setiap \omterm{def:b1:poly:def}{polinomial} bulat yang taknol mempunyai
tinggi semacam itu, dan berakar real paling banyak $\deg P$: sehingga \omterm{pb:b1:logic:1}{bilangan
aljabarnya} membentuk gabungan terbilang (atas $h$) dari \omterm{def:b1:counting:card}{himpunan hingga}, jadi
dapat didaftar sebagai satu barisan tunggal. Seandainya \omterm{pb:b1:logic:1}{bilangan transendennya} juga
dapat didaftar, maka menyelang-selingkan kedua daftarnya akan mendaftar $\R$,
yang bertentangan dengan \cref{pb:b1:reals:1} (pertanyaan 22). Jadi
\omterm{pb:b1:logic:1}{bilangan transenden} membentuk \omterm{def:b1:logic:sets}{himpunan} yang takterbilang. Perbandingannya:
Cantor membuktikan bahwa \emph{kebanyakan} bilangan real bersifat transenden namun tak menunjukkan
satu pun; sedangkan Liouville menunjukkan satu, dengan konstanta yang efektif (Bagian V)
--- jadi keberadaan lewat kelimpahan lawan keberadaan lewat konstruksi.

\textbf{18.} Dari pertanyaan 2, hipotesis lemanya berlaku dengan $s
= 2$, $C = \frac14$. Seandainya $\sqrt2$ bersifat Liouville, maka untuk $n =
3$: $\frac{1}{4q^2} < q^{-3}$ memaksa $q < 4$, sehingga $q \in \{2,
3\}$; lalu hanya berhingga banyak $\frac pq$ dengan $q$ itu yang terletak dalam jarak
$1$ dari $\sqrt2$, masing-masingnya berjarak positif
$\geq \varepsilon_0$ (karena $\sqrt2$ irasional); sehingga memilih $n$ dengan
$2^{-n} < \varepsilon_0$ tak menyisakan $\frac pq$ yang sah sama sekali:
yang bertentangan. Argumen yang sama dengan $\frac{C}{q^d}$ menunjukkan bahwa tak ada
\omterm{pb:b1:logic:1}{bilangan aljabar} yang Liouville --- yaitu pertanyaan 9 dalam pakaian yang
efektif.

\textbf{19.} Di sini $99^2 - 2\cdot70^2 = 9801 - 9800 = 1$. Sehingga
\[
\sqrt2 - \frac{99}{70} =
\frac{2 - (99/70)^2}{\sqrt2 + 99/70}
= \frac{-1}{4900\,\bigl(\sqrt2 + \tfrac{99}{70}\bigr)} ,
\qquad
\Bigl|\sqrt2 - \frac{99}{70}\Bigr|
= \frac{1}{4900 \times 2.8284\dots} \approx 7.2\cdot10^{-5} :
\]
dengan $\frac{99}{70} = 1.414285\dots$ terhadap $\sqrt2 =
1.414213\dots$ --- jadi lima angka yang benar.

\textbf{20.} Uji akar rasional bagi $P = X^3 - 2$: calonnya
$\pm1, \pm2, \pm\frac12$, dan tak satu pun akarnya. Jadi $d = 3$ dan Bagian II
berlaku bagi $x = 2^{1/3} = 1.2599\dots$ Pada $\intcc{x - 1}{x + 1}
\subseteq \intcc{0.25}{2.26}$: $\abs{P'(t)} = 3t^2 \leq 3\,(1 +
2^{1/3})^2 < 3\times(2.26)^2 = 15.32 < 16$, sehingga $M < 16$ dan $C
\geq \frac{1}{16}$:
\[
\Bigl|2^{1/3} - \frac pq\Bigr| \geq \frac{1}{16\,q^3}
\qquad\text{untuk setiap bilangan rasional.}
\]

\textbf{21.} Jika $\bigl|2^{1/3} - \frac pq\bigr| < 10^{-6}$, maka
$\frac{1}{16 q^3} < 10^{-6}$, yakni $q^3 > \frac{10^6}{16} =
62\,500$; dan karena $39^3 = 59\,319 < 62\,500 \leq 64\,000 = 40^3$:
maka $q \geq 40$.

\textbf{22.} Jalankanlah Bagian III dalam basis $2$: $B = \sup_k \sum_{n\leq
k} 2^{-n!}$, $q_k = 2^{k!}$, dan ekor geometrinya (dengan rasio
$\frac12$) memberikan $0 < B - \frac{p_k}{q_k} \leq 2\cdot2^{-(k+1)!}
= 2\,q_k^{-(k+1)} \leq q_k^{-n}$ untuk $k \geq n$. Jadi $B$ bersifat
Liouville, sehingga transenden: jadi tak ada apa pun pada argumennya yang
bersifat desimal.

\textbf{23.} Dengan angka satu pada posisi $3^k$: $q_k = 10^{3^k}$ dan
batas ekornya memberikan $0 < x^\dagger - \frac{p_k}{q_k} <
2\cdot10^{-3^{k+1}} = 2\,q_k^{-3}$ (karena $3^{k+1} = 3\cdot3^k$):
jadi hampiran berorde $3$ yang tak hingga banyaknya. Menurut lema
pertanyaan 15: orde $3 > 1$ menyingkirkan kerasionalannya, dan orde $3 >
2$ menyingkirkan kemungkinannya menjadi irasional kuadratik (yang ketaksamaan
Liouvillenya mempunyai $s = d = 2$). Padahal \omterm{pb:b1:logic:1}{bilangan aljabar} berderajat
$\geq 3$ hanya tertolak pada orde $d \geq 3$: sehingga metode
Liouville tak dapat memisahkan $x^\dagger$ dari yang kubik. Adapun celahnya
ditutup oleh teorema Roth --- bahwa setiap irasional aljabar mempunyai
orde hampiran tepat $2$ --- yaitu hasil abad kedua puluh yang
jauh di luar jilid ini; dan dengan menerimanya, $x^\dagger$ pun
bersifat transenden.

\textbf{24.} Ada paling banyak $(2H+1)^{d+1}$ tupel $(a_0,
\dots, a_d)$ yang entrinya di $\intint{-H}{H}$, dan masing-masing
\omterm{def:b1:poly:def}{polinomial} taknol di antaranya berakar real paling banyak $d$: sehingga paling banyak
$d\,(2H+1)^{d+1}$ \omterm{pb:b1:logic:1}{bilangan aljabar} yang muncul --- yaitu kehinggaan
yang membiarkan pertanyaan 17 mendaftar semuanya.

\textbf{25.} (i) Satu-satunya bahan analisisnya adalah teorema nilai
rata-rata, pada pertanyaan 7, yang mengubah kelenyapan $P(x) =
0$ menjadi tolakan Lipschitz $\abs{P(p/q)} \leq M\abs{x -
p/q}$. (ii) Ketegangannya: bahwa kebulatan mendorong $\abs{P(p/q)}$ naik ke
$q^{-d}$, sedangkan kemulusan menariknya turun ke $M\abs{x - p/q}$ --- sehingga
bilangan rasional yang terlalu dekat dengan $x$ akan tergilas di antara keduanya.
(iii) Liouville mengonstruksi satu \omterm{pb:b1:logic:1}{bilangan transenden} dengan konstanta yang
efektif; sedangkan Cantor menunjukkan bahwa hampir semua bilangan real bersifat transenden
tanpa menamai satu pun: jadi konstruksi lawan \omterm{def:b1:counting:card}{kardinalitas}. (iv) Adapun
soal akhir pekan \cref{ch:b1:integration} membuktikan
keirasionalan $\pi$ lewat apitan yang sama --- yaitu sebuah integral yang
akan menjadi bilangan bulat positif namun terperangkap di $\intoo{0}{1}$ ---
dengan pengintegralan menggantikan pendiferensialan sebagai paruh
analisisnya.

\end{solution}
